% ---
% id: transmitter/decoy/prop:tagging
% kind: prop
% title: "Tagging channel for weights at or below the Poisson weights"
% summary: >-
%   One channel, independent of the label, maps each tagged virtual state to its
%   phase-randomised weak coherent state whenever the declared sector weights
%   lie at or below the Poisson weights. The weights therefore need only a
%   certified inequality, never an exact Poisson weight.
% prominence: main
% status: accepted
% uses: [def:wcp-model, def:tagged-source, transmitter/def:source-link]
% mentions: [rem:no-intensity-split, cor:emitted-decoy]
% assumes: []
% refs:
%   - "tupkary2026decoy, Lemma 9.2 with Eq. (145)"
% provenance: newly-derived
% lean:
% notes: "Newly derived in full: the published Lemma 9.2 is the case of exact Poisson weights with N_ent = N_src, part (iv), and is cited as that special case only. The proposition is an identity among the model states and deliberately assumes no attestation: its assumption closure is empty, and the soundness-registry entry of the tagging kind carries A-WCP beside that closure rather than through it (transmitter/def:source-segment (A1))."
% ---

Recall the virtual source of Definition~\ref{def:tagged-source}\,(T3) and
(T4): $\Avirt$ is the direct sum, over the sectors that some letter weights, of
the quantum sectors $S_{\mathsf{n}}$ for $\mathsf{n}\le N_{\mathrm{ent}}$, the
one-dimensional \tref[polarisation flag]{polarisation flags}
$\mathbb{C}f_{\mathsf{n},a}$ for $N_{\mathrm{ent}}<\mathsf{n}\le N_{\mathrm{src}}$
and the one-dimensional \tref[tail flag]{tail flags} $\mathbb{C}f_{a,\mu}$; the
virtual state $\rvirt_x$ of a letter $x=(a,\mu,z)$ puts
$\tilde p(\mathsf{n}\mid\mu)\dyad{\mathsf{n},a}$ on each quantum sector, the
weight $\tilde p(\mathsf{n}\mid\mu)$ on the polarisation flag $f_{\mathsf{n},a}$
of its own polarisation, and the \tref{tail weight} $\tilde w_\mu$ on its tail
flag $f_{a,\mu}$. The tail flag absorbs whatever weight the declared sectors
leave out, so the construction needs the sector weights only to lie at or below
the \tref[Poisson weight]{Poisson weights}: a declared weight may be any number
certified to lie in $[0,P_\mu(\mathsf{n})]$, and none need equal a Poisson
weight. Lowering $\tilde p(\mathsf{n}\mid\mu)$ moves the
difference into $\tilde w_\mu$, and so into the tail flag of every polarisation
at that intensity; Corollary~\ref{cor:emitted-decoy}\,(v) says what else
changes. Each flag, polarisation or tail, is sent to a state that depends on
the flag alone, so one channel, fixed before a letter is drawn, serves every
letter, which is what a \tref{source link} requires of its channel witness
(Definition~\ref{def:source-link}). That holds for any flag family, the tail
flags indexed by $\mu$ included; Remark~\ref{rem:no-intensity-split} says why
the polarisation flags are nonetheless shared across intensities.

\begin{proposition}\label{prop:tagging}
Let a \tref{weak-coherent-pulse model} be given
(Definition~\ref{def:wcp-model}): labels $x=(a,\mu,z)$ with $z$ the
\tref{round type}, unit polarisation vectors
$\ket{a}$, Poisson weights
$P_\mu(\mathsf{n})=e^{-\mu}\mu^{\mathsf{n}}/\mathsf{n}!$ and model states
$\rho_x=\sum_{\mathsf{n}\ge0}P_\mu(\mathsf{n})\,\dyad{\mathsf{n},a}$ on the
two-mode Fock space. Let a tagged source be declared on it
(Definition~\ref{def:tagged-source}) with cutoffs
$0\le N_{\mathrm{ent}}\le N_{\mathrm{src}}$, virtual states $\rvirt_x$ of
\eqref{eq:rvirt} and declared weights $\tilde p(\mathsf{n}\,|\,\mu)$,
$\mathsf{n}\le N_{\mathrm{src}}$, satisfying the inequality
\eqref{eq:weightineq} of Definition~\ref{def:tagged-source}\,(T2), which
reads
\begin{equation}\label{eq:tag-ineq}
  0\le\tilde p(\mathsf{n}\,|\,\mu)\le P_\mu(\mathsf{n})
  \qquad\text{for every intensity }\mu\text{ and every }
  \mathsf{n}\le N_{\mathrm{src}} ,
\end{equation}
with tail weights
$\tilde w_\mu=1-\sum_{\mathsf{n}\le N_{\mathrm{src}}}\tilde p(\mathsf{n}\,|\,\mu)$.
Then:
\begin{enumerate}[label=(\roman*),leftmargin=2.2em]
\item \emph{(the tail states)} for every polarisation $a$ and every intensity
$\mu$ with $\tilde w_\mu>0$,
\begin{equation}\label{eq:tailstate}
  \vartheta_{a,\mu}\coloneqq\frac{1}{\tilde w_\mu}\Big[
  \sum_{\mathsf{n}\le N_{\mathrm{src}}}
  \big(P_\mu(\mathsf{n})-\tilde p(\mathsf{n}\,|\,\mu)\big)\dyad{\mathsf{n},a}
  +\sum_{\mathsf{n}>N_{\mathrm{src}}}P_\mu(\mathsf{n})\,\dyad{\mathsf{n},a}
  \Big]
\end{equation}
is a state on the Fock space. Where $\tilde w_\mu=0$ the declaration carries no
tail flag for $(a,\mu)$ (Definition~\ref{def:tagged-source}) and no tail state
is defined.
\item \emph{(one channel)} the map
\begin{equation}\label{eq:xitag}
\begin{aligned}
  \Xi_{\mathrm{tag}}(\zeta)\coloneqq{}&
  \sum_{\mathsf{n}\le N_{\mathrm{ent}}}
    V_{\mathsf{n}}\,\Pi^{\mathrm v}_{\mathsf{n}}\,\zeta\,
    \Pi^{\mathrm v}_{\mathsf{n}}V_{\mathsf{n}}^{\dagger}
  +\sum_{N_{\mathrm{ent}}<\mathsf{n}\le N_{\mathrm{src}}}\sum_{a}
    \bra{f_{\mathsf{n},a}}\zeta\ket{f_{\mathsf{n},a}}\,\dyad{\mathsf{n},a}\\
  &+\sum_{(a,\mu):\,\tilde w_\mu>0}
    \bra{f_{a,\mu}}\zeta\ket{f_{a,\mu}}\,\vartheta_{a,\mu} ,
\end{aligned}
\end{equation}
in which $\Pi^{\mathrm v}_{\mathsf{n}}$ projects $\Hs{\Avirt}$ onto the quantum
sector $S_{\mathsf{n}}$, $V_{\mathsf{n}}$ is the isometry that embeds that
sector into the Fock space as $\operatorname{span}\{\ket{\mathsf{n},a}\}_a$,
and every sum runs over the sectors of the declared index set $\mathbb{T}$
only, is a channel from $\Avirt$ to the Fock space, and
\begin{equation}\label{eq:tagid}
  \Xi_{\mathrm{tag}}(\rvirt_x)=\rho_x\qquad\text{for every label }x .
\end{equation}
A quantum sector or a polarisation flag that no letter puts weight on is not a
sector of $\Avirt$ by Definition~\ref{def:tagged-source}\,(T3) and contributes
no term to \eqref{eq:xitag}, exactly as for the tail flags.
The channel $\Xi_{\mathrm{tag}}$ depends on the declared weights and the
cutoffs and on no label.
\item \emph{(the role of the inequality)} \eqref{eq:tagid} holds for every
family of weights satisfying \eqref{eq:tag-ineq}. Where \eqref{eq:tag-ineq}
fails at an intensity $\mu$, its two halves fail differently.
\begin{enumerate}[label=(\alph*),leftmargin=2.2em]
\item If $\tilde p(\mathsf{n}\,|\,\mu)>P_\mu(\mathsf{n})$ for some
$\mathsf{n}\le N_{\mathrm{src}}$ and $\tilde w_\mu>0$, the operator
\eqref{eq:tailstate} is not positive semidefinite, for every polarisation $a$,
and the construction is not available: where a letter $x$ has intensity $\mu$,
the map \eqref{eq:xitag} sends the state $\dyad{f_{a_x,\mu}}$ to
$\vartheta_{a_x,\mu}$ and is not positive.
\item If $\tilde p(\mathsf{n}\,|\,\mu)<0$ for some
$\mathsf{n}\le N_{\mathrm{src}}$ while
$\tilde p(\mathsf{n}'\,|\,\mu)\le P_\mu(\mathsf{n}')$ for every
$\mathsf{n}'\le N_{\mathrm{src}}$, then $\tilde w_\mu>0$ and the operator
\eqref{eq:tailstate} is a state for every polarisation $a$. What fails is the
virtual state: for every letter $x$ of intensity $\mu$ the operator $\rvirt_x$
of \eqref{eq:rvirt} is not positive semidefinite or has trace greater than
one, so it is not a state.
\end{enumerate}
\item \emph{(the published case)} for
$\tilde p(\mathsf{n}\,|\,\mu)=P_\mu(\mathsf{n})$ and
$N_{\mathrm{ent}}=N_{\mathrm{src}}$ the virtual states are those of
\cite[Eq.~(145)]{tupkary2026decoy} and \eqref{eq:tagid} is
\cite[Lemma~9.2]{tupkary2026decoy}.
\end{enumerate}
\end{proposition}

\begin{proof}
\emph{Idea.} The channel keeps each quantum sector, prepares from each
polarisation flag the number state it stands for, and prepares from each tail
flag the normalised remainder of the Poisson distribution, which the
inequality makes a state.

(i) Every coefficient of \eqref{eq:tailstate} is non-negative, the first family
by \eqref{eq:tag-ineq}, and the vectors $\ket{\mathsf{n},a}$, $\mathsf{n}\ge0$,
are orthonormal for a fixed $a$, being number states of the one mode of
polarisation $a$ with the orthogonal mode in its vacuum. So the bracket is a
positive trace-class operator with trace
$\sum_{\mathsf{n}\ge0}P_\mu(\mathsf{n})
-\sum_{\mathsf{n}\le N_{\mathrm{src}}}\tilde p(\mathsf{n}\,|\,\mu)
=1-(1-\tilde w_\mu)=\tilde w_\mu$, and $\vartheta_{a,\mu}$ has unit trace.

(ii) \emph{A channel.} The sector projectors $\Pi^{\mathrm v}_{\mathsf{n}}$,
$\dyad{f_{\mathsf{n},a}}$ and $\dyad{f_{a,\mu}}$ are mutually orthogonal and sum
to $\idon{\Avirt}$ (Definition~\ref{def:tagged-source}). Each of the three sums
of \eqref{eq:xitag} is completely positive: the first is a sum of conjugations
by $V_{\mathsf{n}}\Pi^{\mathrm v}_{\mathsf{n}}$; the second and third are
measure-and-prepare maps, a rank-one projective measurement followed by the
preparation of a fixed state, the states $\vartheta_{a,\mu}$ being states by
(i). For the trace, $V_{\mathsf{n}}$ is an isometry on $S_{\mathsf{n}}$ and
every prepared operator has unit trace, so
$\tr\Xi_{\mathrm{tag}}(\zeta)=\sum_{\mathsf{n}\le N_{\mathrm{ent}}}
\tr[\Pi^{\mathrm v}_{\mathsf{n}}\zeta]
+\sum_{\mathsf{n},a}\bra{f_{\mathsf{n},a}}\zeta\ket{f_{\mathsf{n},a}}
+\sum_{(a,\mu):\,\tilde w_\mu>0}\bra{f_{a,\mu}}\zeta\ket{f_{a,\mu}}
=\tr\zeta$, the projectors summed being all the sector projectors of the
declaration. The input is finite dimensional and every output operator is trace
class, so $\Xi_{\mathrm{tag}}$ is a completely positive trace-preserving map
into the trace class of the Fock space.

\emph{The identity.} Fix $x=(a,\mu,z)$. By Definition~\ref{def:tagged-source},
\begin{equation*}
  \rvirt_x=\sum_{\mathsf{n}\le N_{\mathrm{ent}}}
  \tilde p(\mathsf{n}\,|\,\mu)\,\dyad{\mathsf{n},a}_{S_{\mathsf{n}}}
  +\sum_{N_{\mathrm{ent}}<\mathsf{n}\le N_{\mathrm{src}}}
  \tilde p(\mathsf{n}\,|\,\mu)\,\dyad{f_{\mathsf{n},a}}
  +\tilde w_\mu\,\dyad{f_{a,\mu}} ,
\end{equation*}
which is block diagonal along the sectors. The first sum of \eqref{eq:xitag}
maps its first sum to $\sum_{\mathsf{n}\le N_{\mathrm{ent}}}
\tilde p(\mathsf{n}\,|\,\mu)\dyad{\mathsf{n},a}$, the embedding carrying the
vector $\ket{\mathsf{n},a}$ of $S_{\mathsf{n}}$ to the Fock vector of the same
name. The second maps the flag $f_{\mathsf{n},a}$, which carries weight
$\tilde p(\mathsf{n}\,|\,\mu)$ in $\rvirt_x$ while the flags
$f_{\mathsf{n},a'}$ of the other polarisations $a'\ne a$ carry none, to
$\tilde p(\mathsf{n}\,|\,\mu)\dyad{\mathsf{n},a}$.
The third maps the one tail flag of $\rvirt_x$, where $\tilde w_\mu>0$, to
$\tilde w_\mu\vartheta_{a,\mu}$, which by \eqref{eq:tailstate} is the bracket
there. When $\tilde w_\mu=0$ there is no tail flag and $\rvirt_x$ has no third
term; then \eqref{eq:tag-ineq} and $\sum_{\mathsf{n}}P_\mu(\mathsf{n})=1$ force
$\tilde p(\mathsf{n}\,|\,\mu)=P_\mu(\mathsf{n})$ for
$\mathsf{n}\le N_{\mathrm{src}}$ and $P_\mu(\mathsf{n})=0$ beyond, so the
bracket vanishes as well. Adding the three,
\begin{equation*}
  \Xi_{\mathrm{tag}}(\rvirt_x)
  =\sum_{\mathsf{n}\le N_{\mathrm{src}}}\tilde p(\mathsf{n}\,|\,\mu)
   \dyad{\mathsf{n},a}
  +\sum_{\mathsf{n}\le N_{\mathrm{src}}}
   \big(P_\mu(\mathsf{n})-\tilde p(\mathsf{n}\,|\,\mu)\big)\dyad{\mathsf{n},a}
  +\sum_{\mathsf{n}>N_{\mathrm{src}}}P_\mu(\mathsf{n})\dyad{\mathsf{n},a}
  =\rho_x .
\end{equation*}
Nothing in \eqref{eq:xitag} refers to $x$: the state prepared from a flag, a
polarisation flag indexed by $(\mathsf{n},a)$ or a tail flag indexed by
$(a,\mu)$, is a function of the flag alone.

(iii) The first sentence is (ii), whose only hypothesis on the weights is
\eqref{eq:tag-ineq}.

(a) For every $a$ the vectors $\ket{\mathsf{n}',a}$, $\mathsf{n}'\ge0$, are
orthonormal (proof of (i)), so
$\bra{\mathsf{n},a}\vartheta_{a,\mu}\ket{\mathsf{n},a}
=\big(P_\mu(\mathsf{n})-\tilde p(\mathsf{n}\,|\,\mu)\big)/\tilde w_\mu<0$ and
\eqref{eq:tailstate} is not positive semidefinite. A letter $x$ of intensity
$\mu$ carries the tail flag $f_{a_x,\mu}$, which is a sector of $\mathbb{T}$ by
\eqref{eq:tset} because $\tilde w_\mu>0$. On $\dyad{f_{a_x,\mu}}$ the first two
sums of \eqref{eq:xitag} vanish, the sectors being orthogonal, and the third
gives $\vartheta_{a_x,\mu}$; so \eqref{eq:xitag} maps a positive operator to
one that is not positive semidefinite.

(b) Here
$\sum_{\mathsf{n}'\le N_{\mathrm{src}}}\tilde p(\mathsf{n}'\,|\,\mu)
<\sum_{\mathsf{n}'\le N_{\mathrm{src}}}P_\mu(\mathsf{n}')\le1$, the inequality
being strict at $\mathsf{n}'=\mathsf{n}$, where
$\tilde p(\mathsf{n}\,|\,\mu)<0\le P_\mu(\mathsf{n})$, so $\tilde w_\mu>0$; and
the proof of
(i), which reads only the upper half of \eqref{eq:tag-ineq}, makes
$\vartheta_{a,\mu}$ a state for every $a$. Let $x$ be a letter with
$\mu_x=\mu$, write $a=a_x$, and let $t$ be the sector that carries the term of
$\tilde p(\mathsf{n}\,|\,\mu)$ in \eqref{eq:rvirt}: $t=\mathsf{n}$ if
$\mathsf{n}\le N_{\mathrm{ent}}$ and $t=(\mathsf{n},a)$ otherwise. If
$t\in\mathbb{T}$, the component of $\rvirt_x$ in $S_t$ is
$\tilde p(\mathsf{n}\,|\,\mu)$ times the dyad of the unit vector
$\ket{\mathsf{n},a}$, resp.\ $f_{\mathsf{n},a}$, so the expectation of
$\rvirt_x$ in that vector is $\tilde p(\mathsf{n}\,|\,\mu)<0$ and $\rvirt_x$ is
not positive semidefinite. If $t\notin\mathbb{T}$, \eqref{eq:rvirt} has no term
in $t$. By \eqref{eq:tset} every sector on which $x$ puts a positive weight lies
in $\mathbb{T}$, the tail flag $f_{a,\mu}$ among them since $\tilde w_\mu>0$;
so every term that \eqref{eq:rvirt} omits for $x$ has a weight
$\tilde p(\mathsf{n}'\,|\,\mu)\le0$, the term at $\mathsf{n}$ among them, and by
\eqref{eq:tailweight}
\begin{equation*}
  \tr\rvirt_x
  =1-\sum_{\mathsf{n}'\ \text{omitted}}\tilde p(\mathsf{n}'\,|\,\mu)
  \;\ge\;1-\tilde p(\mathsf{n}\,|\,\mu)\;>\;1 .
\end{equation*}

(iv) With $\tilde p=P_\mu$ and $N_{\mathrm{ent}}=N_{\mathrm{src}}$ the middle
sum of $\rvirt_x$ is empty, the low sectors carry the exact Poisson weights on
$\dyad{\mathsf{n},a}$, and the tail flag carries
$1-\sum_{\mathsf{n}\le N_{\mathrm{src}}}P_\mu(\mathsf{n})$. This is the
displayed state of the cited equation with its photon cutoff equal to
$N_{\mathrm{src}}$ and its flag $\dyad{a,\mu}$ equal to $\dyad{f_{a,\mu}}$; the
cited lemma asserts a channel with \eqref{eq:tagid} for every choice of $a$,
$\mu$ and $z$.
\end{proof}
